\begin{textsample}{POS Dim 2 – mg – Score 27.00 – mg/mg\_em\_146.txt}  \label{ex:f2_pos_193}
7.4 Algumas orientações para formação continuada Para implantar uma nova \textbf{política} pública educacional, além de uma boa formulação, é necessário que todos os atores que compõem as \textbf{redes} de ensino a conheçam, entendam seus propósitos, participem de sua construção e se comprometam com a sua efetiva implementação.
Portanto, não se trata apenas de oferecer o acesso e a leitura de documentos, mas, sim, de garantir espaços e instrumentos que promovam reflexões e formação entre os \textbf{profissionais} da educação sobre a relevância da incorporação em suas práticas cotidianas das mudanças propostas pelo \textbf{Currículo} Referência.
Para que a presente proposta atinja seus objetivos e não se limite a uma ação burocrática, todas as instâncias das \textbf{redes} de ensino de Minas Gerais devem estar envolvidas diretamente em sua implementação, o que exige um processo de formação dos \textbf{profissionais} que atuam nas escolas e, também, dos que trabalham no Órgão Central e nas Superintendências Regionais de Ensino ( SRE ).
É necessário que todos os envolvidos tenham a oportunidade de aprofundar seus conhecimentos sobre a BNCC, as orientações para o Novo Ensino Médio, novas metodologias, novas formas de avaliação, trabalho interdisciplinar e integrado entre professores da mesma área de conhecimento e de diferentes áreas.
Os \textbf{profissionais} da educação acompanham e monitoram as Escolas de Ensino Médio são de fundamental importância, pois mediam as relações entre a Secretaria de Educação e as SRE e os que atuam diretamente com estudantes, podendo ajudar na comunicação e formação, não sendo os únicos agentes nesse processo.
Um \textbf{órgão} estratégico da \textbf{rede} \textbf{estadual} de ensino é a Escola de Formação e Desenvolvimento Profissional de Educadores de Minas Gerais⁴⁵, que trabalhará no processo de formação continuada e pode se dedicar à produção de materiais de apoio que visem a organizar e facilitar a leitura da BNCC e do texto final do \textbf{Currículo}.
Nesse sentido, podem ser produzidos cadernos temáticos ( sujeitos do ensino médio, estratégias avaliativas \textbf{formativas}, integração curricular, gestão democrática, projeto de vida, \textbf{itinerários} \textbf{formativos}, dentre outros temas abordados no \textbf{Currículo} Referência ), vídeos, pílulas explicativas e outros materiais que, para além de apresentar o texto do \textbf{Currículo}, podem apontar estratégias para sua efetiva implementação no ambiente escolar.
Esses materiais podem ser tomados como objeto de estudos nas reuniões das atividades extraclasse, sendo que, em cada encontro, pode -se discutir um módulo/capítulo desse material. ⁴⁵ A Escola de Formação e Desenvolvimento Profissional de Educadores de Minas Gerais foi criada pela Lei Delegada nº 180, de 20 de janeiro de 2011.
Inserida na estrutura da Secretaria de Estado de Educação de Minas Gerais, a Escola de Formação tem como objetivo \textbf{precípuo} coordenar os processos de formação em nível macro, oferecendo suporte logístico, operacional, físico e/ou virtual para realização de \textbf{cursos}, seminários e outras estratégias de formação dos \textbf{profissionais} da educação, em sua dimensão, \textbf{profissional}, cultural e ética.
Fonte : Disponível em.
Acesso em 21 de março de 2021.
Um importante ponto para se começar os estudos é entender as \textbf{legislações} que entraram em vigor a partir de 2017, como a Lei 13.415, de 16 de fevereiro de 2017, que alterou a Lei nº 9.394, de 20 de dezembro de 1996, que estabelece as \textbf{Diretrizes} e Bases da Educação Nacional, e determinou uma alteração na \textbf{carga} \textbf{horária} do Ensino Médio de 800 horas para 1.000 horas anuais até 2022, além de definir uma nova organização curricular, que deve ser formada pelos componentes curriculares da BNCC e pelos \textbf{Itinerários} \textbf{Formativos}, com foco nas áreas do conhecimento e na Formação \textbf{Profissional} e Técnica.
Para compreender as \textbf{Diretrizes} Curriculares Nacionais Gerais para a Educação Profissional e Tecnológica ver a Resolução CNE/CP Nº 1, de 5 de janeiro de 2021.
Todo \textbf{profissional} da educação deve conhecer a BNCC, documento normativo que \textbf{instituiu} os conhecimentos essenciais, as competências gerais para a Educação Básica, habilidades e as aprendizagens pretendidas para todos os estudantes em cada etapa.
Tais orientações nortearam a reelaboração do \textbf{Currículo} Referência.
Outros documentos importantes são a Resolução nº 3, de 21 de novembro de 2018, que \textbf{atualizou} as \textbf{Diretrizes} Curriculares Nacionais para o Ensino Médio ; o \textbf{Guia} de Implementação do Novo Ensino Médio ; os Referenciais Curriculares para a Elaboração de \textbf{Itinerários} \textbf{Formativos}, \textbf{Portaria} nº 1.432, de 28 de dezembro de 2018, em lei ou em formato de livreto, e tantos outros documentos como BNCC Comentada, propostas práticas, documentos e livros elaborados por diferentes instituições.
Ao pesquisar materiais para estudo, devem ser escolhidos os que têm propostas orientadas pelos princípios éticos, políticos e estéticos traçados pelas \textbf{Diretrizes} Curriculares Nacionais da Educação Básica e pelas novas \textbf{Diretrizes} Nacionais para o Ensino Médio e defendem uma educação voltada para a formação humana integral e para a construção de uma sociedade justa, democrática e inclusiva.
A leitura e entendimento da BNCC, do \textbf{Currículo} Referência, das \textbf{Diretrizes}, Decretos, \textbf{Portarias}, dentre outros documentos, são importantes, porque levarão aos Gestores, Analistas, Professores e demais \textbf{profissionais} a compreensão e reflexão em relação às mudanças estruturais e pedagógicas propostas para o Novo Ensino Médio e poderão confrontar ideias e experiências vivenciadas com os problemas encontrados na prática docente e a busca por respostas.
Além do conhecimento das \textbf{legislações}, é imprescindível que os \textbf{profissionais} participem de \textbf{cursos}, como, por exemplo, o \textbf{ofertado} pela Escola de Formação e Desenvolvimento Profissional de Educadores de Minas Gerais, que contempla os estudos da BNCC e os aspectos históricos e processuais da construção do \textbf{Currículo} Referência para a Educação Infantil e o Ensino Fundamental, os eixos estruturadores, as competências gerais, as estratégias de leitura e interpretação das habilidades e a estrutura desse \textbf{Currículo}.
Essa mesma instituição \textbf{ofertará} \textbf{cursos} para que o professor entenda a BNCC para o Ensino Médio e os temas que são primordiais para os \textbf{profissionais} da educação desenvolverem o \textbf{Currículo} Referência nas escolas e nas salas de aula.
Outra possibilidade são os \textbf{cursos} disponíveis na Plataforma do MEC, que são sempre \textbf{atualizados}.
Dentro da \textbf{carga} \textbf{horária} \textbf{semanal} de trabalho correspondente a um cargo de Professor de Educação Básica, há números de horas \textbf{destinadas} às atividades extraclasse, que podem ser momentos importantes de reuniões, discussões, preparação e estudos dos professores dentro do espaço escolar.
Para Gatti ( 2008 ), educação continuada : > [... ] ora se restringe o significado da expressão aos limites de \textbf{cursos} estruturados e formalizados oferecidos após a \textbf{graduação}, ou após ingresso no exercício do magistério, ora ele é tomado de modo amplo e genérico, como compreendendo qualquer tipo de atividade que venha a contribuir para o desempenho \textbf{profissional} – horas de trabalho coletivo na escola, reuniões pedagógicas, trocas cotidianas com os pares, participação na gestão escolar, congressos, seminários, \textbf{cursos} de diversas naturezas e formatos, oferecidos pelas Secretarias de Educação ou outras instituições para pessoal em exercício nos sistemas de ensino, relações \textbf{profissionais} virtuais, processos diversos a distância ( vídeo ou teleconferências, \textbf{cursos} via internet etc. ), grupos de sensibilização \textbf{profissional}, enfim, tudo que possa oferecer ocasião de informação, reflexão, discussão e trocas que favoreçam o aprimoramento \textbf{profissional}, em qualquer de seus ângulos, em qualquer situação ( GATTI, 2008, p. 57 ).
Nessa direção, a formação continuada pode ser feita a partir de atividades coletivas na escola, reuniões pedagógicas, trocas com os pares, \textbf{oficinas}, momentos de reflexão, discussões para o aprimoramento \textbf{profissional}.
Após o primeiro contato com as \textbf{legislações} e \textbf{cursos} introdutórios, sugere -se a segunda fase da formação, que deve acontecer na escola, de forma integrada e contínua.
Esses momentos devem ser organizados ou pelos Gestores, pelos Especialistas ou pelos próprios Professores, para que possam discutir sobre suas práticas pedagógicas, os estudantes, os conhecimentos e dificuldades em relação às novas tecnologias de informação e comunicação e refletir sobre como colocar em prática esse novo \textbf{currículo}, considerando a realidade e o contexto daquela escola.
Um material que pode ser usado para orientar esses momentos de formação na escola são os cadernos do Pacto pelo \textbf{Fortalecimento} do Ensino Médio ( PNEM ), que foi um \textbf{curso} de formação para os \textbf{profissionais} da Educação do Ensino Médio em busca de um ensino mais integrado, inclusivo e participativo, considerando ser essa a proposta do Novo Ensino Médio.
A primeira etapa do PNEM foi composta pela formação comum dos participantes com temas sobre formação integral, integração curricular, organização e gestão democrática, dentre outras.
Na segunda etapa, foram abordados temas sobre organização do trabalho pedagógico no Ensino Médio e de cada área de conhecimento e a integração entre todas elas.
No final deste \textbf{capítulo} há o link para acesso desses materiais.
Reuniões serão úteis para que toda a equipe pedagógica discuta e reflita, principalmente, sobre a prática pedagógica que a escola vem desenvolvendo e as mudanças que ela precisa promover para o atendimento dos objetivos do Novo Ensino Médio.
Pois, para se alcançar as mudanças e melhorias previstas é necessário oferecer uma educação integral aos estudantes de forma integrada, interdisciplinar, com novas metodologias e didáticas e com inserção de pesquisas, planos de intervenção na escola e na sociedade, em espaços dentro e fora da escola.
Essa proposta é possível e acontecerá nas escolas quando mudarmos algumas concepções como : > [... ] hábitos assimilados pela rotina escolar, tais como \textbf{fragmentação} dos espaços/tempos de aprendizagem, objetivação dos procedimentos de ensino, padronização de métodos pedagógicos, agrupamento de alunos por idade ou nível de aprendizagem, \textbf{fragmentação} do \textbf{currículo} por disciplinas, estruturação de \textbf{horários} recortados e rígidos, hierarquização das relações intraescolares, distribuição dos conteúdos escolares por tópicos, determinação de espaços relativamente fixos para cada aluno nas salas de aula, e tantos outros aspectos que marcam a chamada cultura escolar, são construções sociais engendradas desde a Idade Média e reconceitualizadas pela modernidade ( THIESEN, 2011, p. 245 ).
Para mudar essa prática que está incrustada desde a Idade Média, é importante que os \textbf{profissionais} da educação reflitam sobre suas práticas a partir do contexto da escola onde trabalham e da realidade dos estudantes que são \textbf{atendidos}.
Para isso, pensando nas habilidades do século XXI, o maior desafio do professor é se reconhecer como transformador de sua própria realidade por meio de sua prática pedagógica.
Esse poder transformador se torna o processo dialógico entre o ser e o fazer, construindo assim sua autoria permanente frente às necessidades dos jovens estudantes desse século.
Sugestão de materiais : - O \textbf{Guia} de Implementação do Novo Ensino Médio, disponível em http://novoensinomedio.mec.gov.br/resources/downloads/pdf/Guia.pdf - Referenciais Curriculares para a Elaboração de \textbf{Itinerários} \textbf{Formativos}, \textbf{Portaria} nº 1.432, de 28 de dezembro de 2018, ou em formato de livreto, disponível em http://novoensinomedio.mec.gov.br/resources/downloads/pdf/DCEIF.pdf - \textbf{Cursos} disponíveis em http://avamec.mec.gov.br//curso/listar. - Cadernos do Pacto estão disponíveis em : http://www.observatoriodoensinomedio.ufpr.br/pacto-nacional-pelo-fortalecimento-do-ensino-medio/.
Acesso em : 20 fev. 2021.

% matched lemmas: atender, atualizar, capítulo, carga, currículo, curso, destinar, diretriz, estadual, formativo, fortalecimento, fragmentação, graduação, guia, horário, instituir, itinerário, legislação, ofertar, oficina, política, portaria, precípuo, profissional, rede, semanal, órgão
\end{textsample}
